\section{TTLM 中各 TT 核心张量之间的关系}
\label{sec: relationship}

为帮助读者理解 TTLM 中各 TT 核心张量（core）的作用，我们在此给出 TTLM 计算文本 $X = [x^{(1)}, x^{(2)}, \cdots ,x^{(N)}]$ 概率的详细过程。注意，所有中间 TT 核心张量彼此相等：$\tG = \tG^{(2)},...,\tG^{(N-1)}$，且 $\tG^{(1)} = \tG^{(N)}$。

$\vy^{(t)}$（即在时刻 $t$ 给定 $x^{(1:t-1)}$ 时 $x^{(t)}$ 的条件概率）的计算可分为三步。作为第 I 步，假设 $\vf(x^{(1)})$ 是满足 $f(x^{(1)})_1 = 1$ 的 one-hot 向量。TTLM 中 $\tG^{(1)}\vf(x^{(1})$ 的计算如下：

\begin{align}
\tG^{(1)}\vf(x^{(1}) &=
\left[\begin{array}{c}
f\left(x^{(1)}\right)_1 \\
f\left(x^{(1)}\right)_2 \\
\ldots \\
f\left(x^{(1)}\right)_{|V|}
\end{array}\right] \left[\begin{array}{cccc}
\tG_{11}^{(1)} & \tG_{12}^{(1)} & \ldots & \tG_{1 R}^{(1)} \\
\tG_{21}^{(1)} & \tG_{22}^{(1)} & \ldots & \tG_{2 R}^{(1)} \\
\ldots & \ldots & \ldots & \ldots \\
\tG_{|V| 1}^{(1)} & \tG_{|V| 2}^{(1)} & \ldots & \tG_{|V| R}^{(1)}
\end{array}\right] \nonumber \\
&=\left[\tG_{11}^{(1)}, \tG_{12}^{(1)}, \cdots, \tG_{1R}^{(1)}\right]^T \nonumber \\
&=\vh_{\textrm{TTLM}}^{(1)} \nonumber
\end{align}

% \begin{wrapfigure}[13]{r}{0.42\textwidth}
% \small
% \vspace{-18mm}
% \begin{center}
% \includegraphics[scale=0.35]{figure/G_LA.png}
% \vspace{-4mm}
% \label{fig: probability_ttlm}
% \vspace{-8mm}
% \end{center}
% \end{wrapfigure}

作为第 II 步，TTLM 将计算
$\vf(x^{(i)})\tG\vh^{(i-1)}_{\textrm{TTLM}}$，其中 $i \in \{2,3,\cdots,t-1\}$。例如，$\vh_{\textrm{TTLM}}^{(2)}$ 在时刻 $t= 2$ 按式 \ref{eq: ttlm cell full} 计算如下：
\begin{align}
    \vh_{\textrm{TTLM}}^{(2)} = \vf(x^{(2)})^T\tG  \vh^{(1)}_{\textrm{TTLM}} \nonumber
\end{align}

作为第 III 步，TTLM 将输出 $\vy^{(t)}$，如下：

\begin{align}
\tG^{(t)}\vh^{(t-1)}_{\textrm{TTLM}} &=
\left[\begin{array}{cccc}
\tG_{11}^{(t)} & \tG_{12}^{(t)} & \ldots & \tG_{1 R}^{(t)} \\
\tG_{21}^{(t)} & \tG_{22}^{(t)} & \ldots & \tG_{2 R}^{(t)} \\
\ldots & \ldots & \ldots & \ldots \\
\tG_{|V| 1}^{(t)} & \tG_{|V| 2}^{(t)} & \ldots & \tG_{|V| R}^{(t)}
\end{array}\right] \left[\begin{array}{c}
h^{(t-1)}_{{\textrm{TTLM}}_1} \\
h^{(t-1)}_{{\textrm{TTLM}}_2} \\
\ldots \\
h^{(t-1)}_{{\textrm{TTLM}}_R}
\end{array}\right] \nonumber \\
&= \left[\begin{array}{c}
\sum_{i=1}^{R}\tG_{1i}^{(t)} h^{(t-1)}_{{\textrm{TTLM}}_1}  \\
\sum_{i=1}^{R}\tG_{2i}^{(t)} h^{(t-1)}_{{\textrm{TTLM}}_2} \\
\ldots \\
\sum_{i=1}^{R}\tG_{Ri}^{(t)} h^{(t-1)}_{{\textrm{TTLM}}_R}
\end{array}\right] \nonumber
\end{align}

由上述计算可以观察到，$\tG^{(1)}$、$\tG$ 与 $\tG^{(t)}$ 在理论上没有共享参数（尽管为简单起见我们令 $\tG^{(1)} = \tG^{(t)}$）。此外，它们在 TTLM 中扮演的角色各不相同：$\tG^{(1)}$ 可视为词嵌入矩阵；$\tG$ 处理两种来源的信息，即隐藏状态与输入词；$\tG^{(t)}$ 则提取 $\vh^{(t-1)}_{\textrm{TTLM}}$ 中所提供的证据，并在词表上生成一组分数。
